\apendice{Anexo de sostenibilización curricular}

\section{Introducción}

Este anexo recoge una reflexión personal sobre los aspectos de sostenibilidad abordados durante el desarrollo del Trabajo de Fin de Grado. El proyecto consiste en una aplicación web orientada a la gestión de horarios, reservas de aulas, ocupación de espacios y calendario académico dentro del ámbito universitario. Aunque se trata de un trabajo principalmente técnico, su aplicación tiene relación con varios Objetivos de Desarrollo Sostenible (ODS), ya que busca mejorar la organización de recursos, reducir procesos manuales y facilitar una gestión académica más eficiente.

Siguiendo las directrices de sostenibilización curricular de la CRUE \cite{crue2012sostenibilidad}, se analiza a continuación cómo el proyecto se relaciona con la sostenibilidad desde una perspectiva social, ambiental, institucional y técnica.

\section{ODS 4: Educación de calidad}

El proyecto se relaciona directamente con el ODS 4, Educación de calidad, al estar orientado a mejorar la gestión de la actividad académica universitaria. La correcta planificación de horarios, aulas y reservas resulta fundamental para el desarrollo de la docencia, ya que permite organizar los espacios disponibles y evitar conflictos que puedan afectar al alumnado o al profesorado.

La aplicación permite consultar horarios, gestionar reservas periódicas asociadas a la docencia, registrar reservas puntuales y visualizar la ocupación de las aulas. Todo ello contribuye a que la información académica sea más accesible, ordenada y actualizada. De esta forma, se facilita el trabajo de los gestores y se reduce la posibilidad de errores derivados de una gestión manual o dispersa.

Además, la centralización de la información permite que los cambios en el calendario académico, como festivos o cambios de docencia, puedan reflejarse en el sistema de manera más clara. Esto favorece una mejor coordinación entre los distintos agentes implicados en la organización docente.

\section{ODS 9: Industria, innovación e infraestructura}

El proyecto también se alinea con el ODS 9, Industria, innovación e infraestructura, al proponer una solución tecnológica para modernizar un proceso administrativo universitario. La aplicación sustituye parcialmente procedimientos basados en hojas de cálculo y comprobaciones manuales por un sistema web estructurado, con una base de datos centralizada y módulos diferenciados para horarios, reservas, aulas, calendario y usuarios.

Desde el punto de vista técnico, se ha buscado construir una solución mantenible y ampliable. La separación entre \textit{backend} y \textit{frontend}, el uso de una arquitectura modular y la documentación del sistema facilitan que el proyecto pueda evolucionar en el futuro. Esto es importante porque una solución sostenible no solo debe resolver una necesidad inmediata, sino también permitir futuras mejoras sin rehacer todo el sistema.

Asimismo, el proyecto incorpora herramientas de control de versiones y análisis de calidad del código, lo que contribuye a mejorar la mantenibilidad y fiabilidad de la aplicación.

\section{ODS 12: Producción y consumo responsables}

El ODS 12, Producción y consumo responsables, se aborda de forma indirecta mediante la mejora en el uso de los recursos universitarios. Las aulas, laboratorios y espacios docentes son recursos limitados que deben utilizarse de forma eficiente. Una mala planificación puede provocar solapamientos, infrautilización de espacios o necesidad de realizar correcciones manuales posteriores.

La aplicación ayuda a conocer la ocupación real de las aulas y a validar restricciones antes de crear o modificar reservas. Esto permite evitar conflictos y favorece una utilización más responsable de los espacios disponibles. Además, al digitalizar la gestión de horarios y reservas, se reduce la dependencia de documentos impresos o ficheros duplicados, promoviendo una gestión más ordenada de la información.

\section{ODS 13: Acción por el clima}

Aunque el impacto ambiental del proyecto es indirecto, también puede relacionarse con el ODS 13, Acción por el clima. La digitalización de procesos contribuye a reducir el uso de papel y otros recursos asociados a la gestión tradicional de documentos. Además, una planificación más eficiente de aulas y horarios puede servir como base para futuras decisiones relacionadas con el uso energético de los edificios universitarios.

Si en un futuro la aplicación se ampliara con informes de ocupación o análisis de uso de espacios, podría ayudar a identificar aulas infrautilizadas, franjas horarias con baja ocupación o patrones de uso que permitieran optimizar mejor los recursos físicos de la universidad.

\section{Conclusión}

La realización de este Trabajo de Fin de Grado me ha permitido comprender que la sostenibilidad también está presente en el desarrollo de software. Diseñar una aplicación mantenible, segura y útil para mejorar procesos reales forma parte de una visión responsable de la ingeniería informática.

En conclusión, el proyecto contribuye a la sostenibilidad al mejorar la gestión académica, favorecer el uso eficiente de aulas, reducir procesos manuales y ofrecer una base tecnológica preparada para futuras ampliaciones. Aunque su impacto ambiental no es directo, sí aporta valor desde una perspectiva social, institucional y técnica dentro del contexto universitario.

